\documentclass[12pt, a4paper]{article}

% --- PREAMBUŁA (ustawienia dokumentu) ---
\usepackage[utf8]{inputenc}       % Kodowanie znaków
\usepackage[T1]{fontenc}          % Polskie znaki (ą, ć, ę, ł, ń, ó, ś, ź, ż)
\usepackage[polish]{babel}        % Polskie reguły dzielenia wyrazów i nazwy (np. "Spis treści")
\usepackage{amsmath}              % Pakiety matematyczne (niezbędne do macierzy)
\usepackage{amsfonts}             % Czcionki matematyczne
\usepackage{amssymb}              % Symbole matematyczne (np. nieskończoność)
\usepackage{geometry}             % Ustawienia marginesów
\usepackage{csvsimple}
\usepackage{float}
\usepackage{graphicx}
\usepackage{siunitx}
\usepackage[utf8]{inputenc}
\usepackage[T1]{fontenc}
\usepackage{booktabs} % Wymagane przez pandas.to_latex()
\usepackage{textcomp} % Dla symbolu mikrona
\geometry{
	a4paper,
	left=2.5cm,
	right=2.5cm,
	top=2.5cm,
	bottom=2.5cm
}
\usepackage{indentfirst}          % Wcięcia w pierwszych akapitach
\usepackage{hyperref}             % Aktywne odnośniki w dokumencie
\usepackage{setspace}             % Interlinia
\setstretch{1.35}                 % ZWIĘKSZONA INTERLINIA (ok. 1.5 linii)
\setlength{\parskip}{0.5em}       % Dodatkowy odstęp między akapitami

% --- DANE DO STRONY TYTUŁOWEJ ---
\title{\textbf{Sprawozdanie z projektu nr 1} \\ \large Asymetryczny Problem Komiwojażera (ATSP)
\\ \large Projektowanie Efektywnych algorytmów}
\author{Karol Gąsior / 280913}
\date{\today} 

% --- POCZĄTEK WŁAŚCIWEGO DOKUMENTU ---
\begin{document}
	
	\maketitle % Generuje tytuł, autora i datę
	\newpage
	\tableofcontents % Automatyczny spis treści
	\newpage
	
\section{Wstęp teoretyczny}

\subsection{Algorytm podziału i ograniczeń (Branch and Bound)}
W ramach drugiego etapu projektu zaimplementowano algorytm podziału i ograniczeń (ang. \textit{Branch and Bound}, B\&B) w celu rozwiązania asymetrycznego problemu komiwojażera (ATSP)[cite: 2, 3]. Metoda ta jest dokładnym algorytmem optymalizacyjnym, który w przeciwieństwie do przeglądu zupełnego nie sprawdza każdej możliwej permutacji w sposób naiwny. 

Działanie algorytmu opiera się na budowaniu drzewa przestrzeni rozwiązań. W każdym węźle drzewa obliczane jest tzw. dolne ograniczenie (ang. \textit{Lower Bound}, LB) kosztu dla wszystkich ścieżek wychodzących z danego węzła. Jeżeli wartość LB danej gałęzi przekracza lub jest równa kosztowi najlepszego dotychczas znalezionego rozwiązania (górne ograniczenie), cała gałąź jest odrzucana (odcinana). Zgodnie z założeniami projektowymi, algorytm został wyposażony w trzy metody przeglądania tak zbudowanej przestrzeni: \textit{breadth-first-search} (przeszukiwanie wszerz), \textit{depth-first-search} (przeszukiwanie w głąb) oraz \textit{lowest-cost} (najlepszy najpierw)[cite: 6].

% ---------------------------------------------------------------------------------------------------------------------

\section{Przykład praktyczny algorytmu B\&B (Depth-First Search) dla $N=3$}

Poniżej przedstawiono działanie algorytmu krok po kroku dla przykładowej instancji problemu o małej wartości $N$[cite: 14]. Przyjęto strategię przeszukiwania w głąb (DFS), a dolne ograniczenie wyznaczane jest metodą Little'a. Niech początkowa asymetryczna macierz kosztów $M$ wynosi:

$$ M = \begin{bmatrix} \infty & 10 & 15 \\ 5 & \infty & 9 \\ 6 & 12 & \infty \end{bmatrix} $$

\textbf{Krok 1: Węzeł początkowy (Root) i wyznaczenie LB}
\begin{itemize}
	\item \textbf{Redukcja wierszy:} W każdym wierszu znajdujemy wartość minimalną i odejmujemy ją od pozostałych elementów wiersza. 
	Minimum wiersza 0: $10$. Po odjęciu: $[\infty, 0, 5]$.
	Minimum wiersza 1: $5$. Po odjęciu: $[0, \infty, 4]$.
	Minimum wiersza 2: $6$. Po odjęciu: $[0, 6, \infty]$.
	Suma redukcji wierszy: $10 + 5 + 6 = 21$.
	\item \textbf{Redukcja kolumn:} Analogicznie odejmujemy minima z kolumn (dla zredukowanej już macierzy).
	Minimum kolumny 0: $0$.
	Minimum kolumny 1: $0$.
	Minimum kolumny 2: $4$. Po odjęciu: kolumna 2 to $[1, 0, \infty]$.
	Suma redukcji kolumn: $0 + 0 + 4 = 4$.
	\item \textbf{Macierz zredukowana węzła początkowego $M'$:}
	$$ M' = \begin{bmatrix} \infty & 0 & 1 \\ 0 & \infty & 0 \\ 0 & 6 & \infty \end{bmatrix} $$
	\item \textbf{Koszt węzła (LB):} $21 + 4 = 25$. Węzeł trafia na \texttt{Stack}.
\end{itemize}

\textbf{Krok 2: Rozgałęzienie (Branching) z wierzchołka 0}
Pobieramy węzeł ze stosu. Możemy udać się do miasta 1 lub 2. Zgodnie z metodą DFS, badamy obydwie opcje i umieszczamy je na stosie.

\textit{Węzeł Dziecko A (ścieżka 0 $\rightarrow$ 1):}
\begin{itemize}
	\item Kopiujemy macierz $M'$. Blokujemy wiersz 0 i kolumnę 1 (ustawiamy na $\infty$). Blokujemy powrót (ustawiamy $M'[1][0] = \infty$).
	\item Macierz przed ponowną redukcją:
	$$ M_A = \begin{bmatrix} \infty & \infty & \infty \\ \infty & \infty & 0 \\ 0 & \infty & \infty \end{bmatrix} $$
	\item Macierz jest już w pełni zredukowana (w każdym niezablokowanym wierszu i kolumnie jest zero). Koszt redukcji $= 0$.
	\item Nowe $LB_A$ = $LB_{poprzednie}$ (25) + Koszt przejścia $M'[0][1]$ (0) + Koszt redukcji (0) = 25.
\end{itemize}

\textit{Węzeł Dziecko B (ścieżka 0 $\rightarrow$ 2):}
\begin{itemize}
	\item Blokujemy wiersz 0, kolumnę 2 oraz powrót $M'[2][0]$.
	\item Macierz przed redukcją:
	$$ M_B = \begin{bmatrix} \infty & \infty & \infty \\ 0 & \infty & \infty \\ \infty & 6 & \infty \end{bmatrix} $$
	\item Redukcja wiersza 2 (minimum to 6). Po odjęciu wiersz 2 to $[\infty, 0, \infty]$. Suma redukcji $= 6$.
	\item Nowe $LB_B$ = $LB_{poprzednie}$ (25) + Koszt przejścia $M'[0][2]$ (1) + Koszt redukcji (6) = 32.
\end{itemize}
\textbf{Zarządzanie stosem:} Na stos (LIFO) odkładamy Węzeł B, a następnie Węzeł A.

\textbf{Krok 3: Eksploracja DFS (Węzeł A)}
\begin{itemize}
	\item Ściągamy ze stosu Węzeł A (ścieżka: 0 $\rightarrow$ 1, LB = 25). Ponieważ nie odwiedzono jeszcze wszystkich miast, generujemy jedyne możliwe dziecko: ścieżkę 0 $\rightarrow$ 1 $\rightarrow$ 2.
	\item Macierz zostaje całkowicie wyzerowana/zablokowana, koszt nie rośnie. Kompletna trasa to 0 $\rightarrow$ 1 $\rightarrow$ 2 $\rightarrow$ 0.
	\item Ustawiamy najlepszy dotychczasowy koszt rozwiązania (\textit{upper bound}) na 25.
\end{itemize}

\textbf{Krok 4: Próba eksploracji Węzła B i odcięcie (Pruning)}
\begin{itemize}
	\item Ściągamy ze stosu Węzeł B (ścieżka 0 $\rightarrow$ 2, LB = 32).
	\item Ponieważ $LB_B$ (32) jest większe niż dotychczas znalezione najlepsze rozwiązanie (25), algorytm natychmiast odrzuca tę gałąź bez dalszej eksploracji. Stos jest pusty, algorytm kończy działanie.
\end{itemize}

% ---------------------------------------------------------------------------------------------------------------------

\section{Opis implementacji algorytmu i zastosowanych struktur}

Implementacji dokonano w języku C++ z zachowaniem obiektowego paradygmatu programowania[cite: 10]. Mając na uwadze optymalizację pamięciową, wszystkie używane struktury danych są alokowane dynamicznie i zależą wyłącznie od aktualnie wprowadzonego rozmiaru problemu[cite: 5].

\subsection{Kluczowe struktury danych}
\begin{itemize}
	\item \textbf{Struktura \texttt{Node}:} Fundamentalna klasa reprezentująca węzeł w drzewie przestrzeni rozwiązań. Przechowuje zredukowaną macierz kosztów przejść (reprezentującą stan problemu w danym punkcie), listę odwiedzonych miast (wektor ścieżki), indeks obecnego miasta, głębokość w drzewie oraz skalkulowaną dla tego węzła wartość dolnego ograniczenia (\textit{Lower Bound}). Przechowywanie instancji macierzy bezpośrednio w węźle ułatwia implementację metod DFS/BFS i izoluje stany między rozgałęzieniami.
	\item \textbf{Struktura \texttt{Stack}:} Dynamiczna struktura danych typu LIFO (\textit{Last-In, First-Out}). Jest kluczowym elementem implementacji przeglądania przestrzeni metodą \textit{depth-first-search}[cite: 6]. Zastosowanie stosu pozwala na sztuczne wywołania rekurencyjne (iteracyjne zagłębianie się w drzewo) bez ryzyka przepełnienia stosu systemowego (\textit{Stack Overflow}) dla instancji o większym wymiarze $N$. W przypadku strategii \textit{breadth-first-search} zamiast stosu wykorzystywana jest struktura kolejki (\textit{Queue}), a dla \textit{lowest-cost} - kolejka priorytetowa (\textit{Priority Queue}).
\end{itemize}

\subsection{Funkcja obliczająca ograniczenia dla metody B\&B (Algorytm Little'a)}
Funkcja wyznaczająca dolne ograniczenie w każdym węźle opiera się na algorytmie redukcji macierzy (metoda Little'a)[cite: 15]. Procedura ta składa się z następujących etapów:
\begin{enumerate}
	\item \textbf{Zabezpieczenie przed cyklami podrzędnymi:} Przed rozpoczęciem redukcji dla tworzonego węzła-dziecka, ścieżka z miasta $i$ (rodzic) do $j$ (dziecko) jest usuwana z dostępnych opcji. Cały wiersz $i$ oraz kolumna $j$ w macierzy są wypełniane wartościami $\infty$. Dodatkowo, aby zapobiec przedwczesnemu zamknięciu cyklu, krawędź powrotna z $j$ do $i$ również otrzymuje koszt $\infty$.
	\item \textbf{Redukcja wierszowa:} Algorytm iteruje po wszystkich niezablokowanych wierszach. Dla każdego znajduje element o najmniejszej wartości i odejmuje go od wszystkich elementów w tym wierszu. Zliczana jest suma odjętych minimów.
	\item \textbf{Redukcja kolumnowa:} Proces jest powtarzany dla kolumn nowej macierzy. Minima z niezablokowanych kolumn są odejmowane, a ich wartości powiększają sumę całkowitą.
	\item \textbf{Ewaluacja kosztu:} Ostateczna wartość ograniczenia dolnego (\textit{Lower Bound}) dla nowego węzła obliczana jest jako suma trzech elementów: $LB$ węzła rodzica, kosztu krawędzi $(i, j)$ z pierwotnej macierzy rodzica oraz całkowitego kosztu obecnej redukcji wierszowo-kolumnowej.
\end{enumerate}	

	
	
\end{document}